\section{Future Directions}
\label{sec:future}

We identify five directions likely to change how spatial resolution enhancement is approached computationally.

\textbf{1. Foundation models for spatial omics.}
Large-scale pre-trained models---Nicheformer \citep{nicheformer2024}, Novae, OmiCLIP, and spatial extensions of scGPT \citep{cui2024scgpt} and CellPLM---are learning transferable representations of cellular identity and spatial context from atlas-scale data. These models could reduce dependence on per-dataset references by providing universal cell-type embeddings that generalise across tissues and platforms. Whether representations learned primarily from dissociated single-cell data transfer to the spatial domain, where neighbourhood context carries biological meaning absent in suspension, remains an open question.

Among these, Nicheformer explicitly encodes spatial neighbourhood composition during pre-training, making it the most directly relevant to deconvolution and reconstruction tasks. scGPT, by contrast, was designed for dissociated cells and requires fine-tuning on spatial data to capture tissue architecture. The relative merits of spatial-native versus adapted foundation models have not been systematically compared, and doing so will require the standardised benchmarks discussed below.

\textbf{2. The Visium HD era and sub-cellular sequencing.}
Visium HD's 2~\textmu m bins approach single-cell resolution for many tissue types, potentially rendering traditional spot-level deconvolution unnecessary. Yet computational challenges persist: at 2~\textmu m, individual bins capture very few UMIs, requiring aggregation or imputation to achieve reliable expression estimates. The deconvolution problem is thus reframed rather than eliminated---from ``which cell types are in this spot?'' to ``how should we aggregate sparse bins into biologically meaningful units?'' Furthermore, the vast installed base of standard Visium data (55~\textmu m spots) ensures that computational resolution enhancement will remain relevant for re-analysis of existing datasets for years to come.

\textbf{3. Reference-free and self-supervised methods.}
Current reference-free approaches (STdeconvolve, BayesTME) are limited to Stage~I and produce coarse outputs. Self-supervised learning---training on the spatial data itself using spatial consistency, image features, or multi-scale structure as supervision signals---offers a path toward reference-free methods that achieve Stage~II or even Stage~III resolution. Contrastive learning frameworks that enforce consistency between a spot's expression and its histological appearance are a promising direction.

\textbf{4. Standardised benchmarks and community evaluation.}
The field urgently needs shared benchmark resources analogous to the Open Problems in Single-Cell Analysis initiative. Key requirements include: (i) multi-platform datasets with matched sequencing-based and imaging-based ST on the same tissue; (ii) agreed-upon metrics that capture both accuracy and calibration; (iii) standardised preprocessing to eliminate confounding differences between evaluation studies; and (iv) living leaderboards that track method performance as new tools are published.

\textbf{5. Multi-omics spatial integration.}
Emerging technologies measure not only transcriptomes but also proteins (CODEX, IMC), chromatin accessibility (spatial ATAC-seq), and metabolites in spatial context. Integrating these modalities computationally---for example, using spatial proteomics to anchor transcriptomic deconvolution, or using chromatin state to inform cell-type identity---will require new multi-modal frameworks that go beyond the expression-centric methods reviewed here. The challenge is not merely technical but also conceptual: defining what ``single-cell resolution'' means when different modalities have different native resolutions and measurement characteristics.

\textbf{6. Computational cost and accessibility.}
Many recent methods (SpatialScope, CellMap, graph-based approaches) require GPU hardware and hours of computation for a single Visium slide. As spatial transcriptomics becomes routine in clinical settings, methods must scale to hospital-grade hardware and turnaround times. Lightweight approximations, pre-computed reference embeddings, and cloud-based pipelines will be necessary to make resolution enhancement accessible beyond well-resourced computational biology groups.
